\begin{tcolorbox}[colback=red!3!white,colframe=red!50!black,boxrule=0.8pt,title={Korrekturhinweise},fonttitle=\bfseries]
\textbf{Korrekturen aus der Rechenprüfung (30.~Sept.~2026).} Die folgenden Stellen sind im Text korrigiert bzw. vermerkt; jede trägt dort einen datierten Hinweis mit dem vorherigen Wert:
\begin{itemize}
\item Zeuge~A: Die Weg-2-Formel $E_0^2 = 4\sqrt2\,m_\mu\,\xi^{-p}$ ist aus A130 entfernt (nicht hergeleitet; sie setzt $m_e/\text{MeV} = 4\sqrt2\,K\,\xi^{q}$, ist also einheitenabhängig). Zeuge~A ist jetzt das von $\alpha$ verlangte $K_\text{frak} = 0{,}98650$; Residuen $1{,}7\cdot10^{-4}$ (additiv) und $2{,}6\cdot10^{-4}$ (multiplikativ) -- keine Formaussage (vorher $7{,}5{:}1$ additiv mit $p=-(2-\sqrt3)$).
\item Abschnitt \enquote{Nachrechnung der p-Identität} entfernt (Vermerk); P-315-1 entfällt, P-315-2 neu gefasst.
\item Buchungsstand: Form additiv \enquote{zweifach} $\to$ \enquote{einfach} bedingt-bestätigt (nur Zeuge~B).
\item P-315-3 (1.~Okt.~2026): $n$-Weg $101{,}3 \to 101{,}2$, angeglichen an die bereits korrigierte Stelle im Abschnitt zur $n$-Sprache.
\end{itemize}
\end{tcolorbox}

% 315_Kfrak_Form_De_ch.tex -- Kapitelinhalt Dok. 315
% Die Form von K_frak: additiv oder multiplikativ?
% Mit dem Kontrollfall der Eulerschen Musikspirale (5/7-Limit).

\section{Worum es geht}

Der fraktale Korrekturfaktor tritt im Korpus in der additiven Form auf,
\[
  K_\text{frak}^{\,\text{add}} \;=\; 1 - 100\,\xi \;=\; \frac{74}{75}
  \;=\; 0{,}98\overline{6},
\]
w\"ahrend eine multiplikativ-kumulative Form denkbar w\"are:
\[
  K_\text{frak}^{\,\text{mul}} \;=\; (1-\xi)^{100} \;=\; 0{,}9867543\ldots
\]
Der relative Abstand der beiden Formen betr\"agt $8{,}88\cdot10^{-5}$. Der
Korpus best\"atigt den \emph{Wert} $74/75$ mehrfach unabh\"angig (A040, A130,
A270); ob damit auch die \emph{Form} -- additiv gegen multiplikativ --
diskriminiert ist, war bislang ungepr\"uft. Dieses Dokument holt die Pr\"ufung
nach und beantwortet drei Fragen: Welche Korpus-Stelle kann die Formen
aufl\"osen? Was ergibt die Nachrechnung der dabei auftretenden offenen
Identit\"aten? Und -- als Strukturfrage -- \emph{warum} sollte die eine Form
die richtige sein?

Als Kontrollfall dient die Eulersche Musikspirale. Sie ist das klassische
Beispiel eines Zyklus, der \emph{nie} schlie\ss t, und liefert damit die
Gegenprobe zum $\xi$-Zyklus, der nach $75$ Uml\"aufen \emph{exakt} schlie\ss t.
Der Vergleich der beiden F\"alle stellt am Ende die begriffliche Achse, an der
die Formfrage h\"angt: eingerollte gegen ausgerollte Beschreibung.

\textbf{Vorab-Kl\"arung (Kategorienfehler vermeiden):} Der Wert $74/75$ ist per
A040-Herleitung \textbf{[K]}; die Frage dieses Dokuments betrifft
ausschlie\ss lich, ob unabh\"angige Zeugen die additive gegen die
multiplikative Form \emph{messen} k\"onnen -- also den Status der Form als
gemessene statt nur hergeleitete Aussage.

\section{Kontrollfall I: die Eulersche Spirale schlie\ss t nie}

\subsection{Das Unm\"oglichkeitsargument}

\textbf{[Q]} Eulers Tonnetz (\emph{Tentamen novae theoriae musicae}, 1739)
spannt die Intervalle der reinen Stimmung als Gitter aus Quinten ($3/2$) und
gro\ss en Terzen ($5/4$) auf; die Oktave ($2$) ist herausdividiert. Die Frage
\enquote{nach wie vielen Wiederholungen schlie\ss t die Spirale?} hat eine
exakte Antwort: \emph{nie}. Ein Schluss verlangte
\[
  2^{x}\cdot 3^{y}\cdot 5^{z} = 1
  \qquad\text{mit ganzzahligen } (x,y,z)\neq(0,0,0),
\]
was der Eindeutigkeit der Primfaktorzerlegung widerspricht. \"Aquivalent:
$\log 2$, $\log 3$, $\log 5$ sind \"uber $\mathbb{Q}$ linear unabh\"angig. Die
Bahn liegt dicht auf dem Oktavkreis, trifft ihren Ausgangspunkt aber nie.

\subsection{Beinahe-Schl\"usse im 5-Limit}

\textbf{[Q]} Was es stattdessen gibt, sind Beinahe-Schl\"usse, deren
Fehlbetrag jeweils ein benanntes Komma ist:

\begin{center}
\begin{tabular}{llll}
\toprule
Weg auf dem Tonnetz & Restfehler (Komma) & Verh\"altnis & Cent \\
\midrule
8 Quinten $+$ 1 gr.~Terz vs.\ 5 Oktaven & \textbf{Schisma} & $32805/32768$ & $\approx 1{,}95$ \\
4 Quinten vs.\ Terz $+$ 2 Oktaven & syntonisches Komma & $81/80$ & $\approx 21{,}51$ \\
12 Quinten vs.\ 7 Oktaven & pythagoreisches Komma & $531441/524288$ & $\approx 23{,}46$ \\
3 gr.~Terzen vs.\ 1 Oktave & kleine Di\"esis & $128/125$ & $\approx 41{,}06$ \\
\bottomrule
\end{tabular}
\end{center}

Der praktisch beste Schluss im 5-Limit ist das Schisma ($\approx 2$ Cent,
unter der H\"orschwelle); der bekannteste ist der Quintenzirkel nach $12$
Quinten mit dem pythagoreischen Komma als Rest.

\subsection{Beinahe-Schl\"usse im 7-Limit (Skriptergebnis)}

\textbf{[K]} Nimmt man die Naturseptime $7/4$ als dritte Achse hinzu
(Gitter $\mathbb{Z}^3$ aus $3,5,7$), gilt dasselbe
Unm\"oglichkeitsargument. Die systematische Suche
(\texttt{euler\_spirale\_7limit.py}, exakte Bruchrechnung, Suchraum
$|b|\le 12$ Quinten, $|c|\le 6$ Terzen, $|d|\le 4$ Septimen; eine Assertion
best\"atigt mit, dass nie ein exakter Schluss auftritt) liefert
Beinahe-Schl\"usse, die eine Gr\"o\ss enordnung besser sind als alles im
5-Limit:

\begin{center}
\begin{tabular}{llll}
\toprule
Cent & Komma & Weg $3^{b}\,5^{c}\,7^{d}$ & Name \\
\midrule
$0{,}325$ & $250047/250000$ & $b=-6,\ c=6,\ d=-3$ & Landscape-Komma \\
$0{,}396$ & $4375/4374$ & $b=-7,\ c=4,\ d=1$ & \textbf{Ragisma} \\
$0{,}721$ & $2401/2400$ & $b=-1,\ c=-2,\ d=4$ & Breedsma \\
$1{,}954$ & $32805/32768$ & $b=-8,\ c=-1,\ d=0$ & Schisma (5-Limit, Vergleich) \\
$5{,}758$ & $5120/5103$ & $b=-6,\ c=1,\ d=-1$ & Hemifamity \\
$7{,}712$ & $225/224$ & $b=-2,\ c=-2,\ d=1$ & septimales Kleisma \\
\bottomrule
\end{tabular}
\end{center}

Das Ragisma ist der einfachste Sub-Cent-Schluss ($7$ Quinten abw\"arts,
$4$ Terzen und $1$ Septime aufw\"arts verfehlen den Ausgangspunkt um
$0{,}4$ Cent); das gr\"obere 7-Limit-Gegenst\"uck zum syntonischen Komma ist
das Archytas-Komma $64/63$ ($\approx 27$ Cent). Mehr Primachsen machen das
Gitter dichter -- die Spirale kommt sich selbst schneller nahe, schlie\ss t
aber prinzipiell nie.

\subsection{Wie man die Spirale doch schlie\ss t: Temperierung}

\textbf{[B]} Ein Schluss entsteht erst, wenn man Kommas \emph{wegtemperiert},
also k\"unstlich zu $1$ erkl\"art. Die zw\"olfstufige gleichstufige Temperatur
setzt die Quinte per Dekret auf exakt $7/12$ der Oktave -- sie
\emph{rationalisiert} die Schrittweite, und erst dadurch schlie\ss t der
Quintenzirkel. Der geschlossene Zirkel ist keine Eigenschaft der reinen
Intervalle, sondern ein Artefakt der Rationalisierung. Dieser Satz wird in
Abschnitt~\ref{sec:eingerollt} tragend.

\section{Kontrollfall II: der Xi-Zyklus schlie\ss t exakt}

\textbf{[B]} Der $\xi$-Zyklus steht auf der anderen Seite derselben
Dichotomie. Die Schrittweite ist
\[
  100\,\xi = \frac{400}{30000} = \frac{1}{75}
  \qquad\text{-- exakt rational,}
\]
weil $\xi = 4/30000$ eine rationale Setzung ist (A020). Ein Schritt von
$1/75$ eines Zyklus schlie\ss t nach genau $75$ Wiederholungen ohne Rest;
$K_\text{frak} = 1 - 1/75 = 74/75$ ist der geschlossene Weg minus ein
Segment (A040, A050). Das ist die allgemeine Torus-Dichotomie: rationale
Windungszahl $\Rightarrow$ geschlossene Bahn; irrationale $\Rightarrow$
dichte, nie schlie\ss ende Bahn.

\begin{center}
\begin{tabular}{lll}
\toprule
 & Musikspirale (5/7-Limit) & $\xi$-Zyklus (FFGFT) \\
\midrule
Schrittweite & $\log_2(3/2),\ \log_2(5/4),\ldots$ irrational & $1/75$ rational \\
Grund & Primzahlen mult.\ unabh\"angig & $\xi$ rational gesetzt \\
Verhalten & dicht, trifft sich nie & schlie\ss t exakt nach 75 \\
\enquote{Fast-Schluss} & Kommas (Schisma, Ragisma, \ldots) & entf\"allt \\
Verifikation & nur numerisch (Cent) & exakte Bruchrechnung \\
\bottomrule
\end{tabular}
\end{center}

Die Schlie\ss ung nach $75$ ist also kein numerischer Zufall, sondern
direkte Folge der Rationalit\"at von $\xi$ -- dieselbe Rationalisierung, die
in der Musik erst die Temperierung nachtr\"aglich erzwingt, ist im Rahmen
von Anfang an gesetzt.

\section{Die Formfrage}
\label{sec:formfrage}

\textbf{[B]} Die beiden Formen sind bis zur ersten Ordnung identisch; ihr
gesamter Unterschied ist der zweite Binomialterm:
\[
  (1-\xi)^{100} = 1 - 100\xi + \binom{100}{2}\xi^2 - \ldots,
  \qquad
  K^{\text{mul}} - K^{\text{add}} = 4950\,\xi^2 - O(\xi^3)
  = 8{,}76\cdot10^{-5}.
\]
Numerisch verifiziert: der Binomialterm $4950\xi^2 = 8{,}8\cdot10^{-5}$
deckt $99{,}6\,\%$ des vollen Abstands; der Rest sind h\"ohere Terme. Eine
Korpus-Stelle kann die Form nur entscheiden, wenn ihre eigene Unsicherheit
kleiner ist als dieser Abstand -- bei Potenz-Stellen: kleiner als der
entsprechend verst\"arkte Abstand.

\section{Drei Zeugen}
\label{sec:zeugen}

\subsection{Zeuge A: die $\alpha$-Stelle (A130)}

\textbf{[K]} A130 führt zur Skala $E_0$ zwei eng verwandte Wege: Weg~1 als
nacktes geometrisches Mittel $E_0 = \sqrt{m_e m_\mu} = 7{,}348$~MeV, Weg~2
als dasselbe Mittel mit fraktaler Korrektur,
$E_0 = \sqrt{m_e m_\mu/K_\text{frak}} = 7{,}397$~MeV. Sie unterscheiden sich
nur um $K_\text{frak}$. Das gemessene $\alpha$ legt diesen Faktor fest,
ohne ihn zu benutzen:
\[
  K_\text{mess} = \xi\,m_e m_\mu\,\alpha^{-1} = 0{,}98650 .
\]

\textbf{[K]} Befund: Das Residuum gegen die additive Form beträgt
$1{,}7\cdot10^{-4}$, gegen die multiplikative $2{,}6\cdot10^{-4}$. Beide
liegen weit über dem Formabstand $8{,}9\cdot10^{-5}$; die Stelle hat einen
eigenen Rest, der größer ist als der Abstand der Formen. Zeuge~A bestätigt
den \emph{Wert} $74/75$ auf $1{,}7\cdot10^{-4}$, macht aber \emph{keine
Formaussage}. \textit{(Korrigiert am 30.~Sept.~2026: vorher war Weg~2 die Konstruktion $E_0^2 = 4\sqrt2\,m_\mu\,\xi^{-p}$ ohne $m_e$; mit dem Zahlenliteral $p=-0{,}2679$ ergab sich $K_\text{mess} = 0{,}986220$, mit $p=-(2-\sqrt3)$ exakt $0{,}9866533$ und ein Urteil $7{,}5{:}1$ additiv. Diese Formel ist aus A130 als zweiter Weg entfernt: $p$ ist im Korpus nicht hergeleitet (so auch dieses Dokument), und da sich $m_\mu$ kürzt, sagt sie nur $m_e/\text{MeV} = 4\sqrt2\,K\,\xi^{q}$ -- eine reine Zahl gleich einer Energie in der willkürlichen Einheit MeV. Ihre Nähe zu $2-\sqrt3$ hängt damit an der Einheit, wie die $e^2$-Koinzidenz in A130.)}

\textbf{[K] Registerlage (R72, eingearbeitet 3.~August 2026):} Das
Register bucht die beiden $E_0$-Werte als \emph{denselben Wert nackt und
fraktal korrigiert} ($7{,}348$ nackt, $7{,}398$ $\alpha$-verankert;
Verhältnis $1/\sqrt{K_\text{frak}}$ nach der A040-Regel). Aus dem
Verhältnis folgt eine Bestimmung des \emph{Exponenten}:
$n = 101{,}2$, neben $n = 100\pm2$ (R67) und $n = 100{,}27$
(Sektorleiter A270). In dieser $n$-Sprache lautet die Formfrage: additiv
$K = 1-n\xi$ gegen multiplikativ $(1-\xi)^{n}$ unterscheiden sich um
$\Delta n_\text{eff} = n(n-1)\xi/2 \approx 0{,}66$ -- keine der drei
$n$-Bestimmungen erreicht diese Schärfe derzeit. \textit{(Korrigiert am 30.~Sept.~2026: vorher $n = 101{,}3$ und \enquote{Was bleibt, ist die exakte $q^{*}$-Bestimmung \ldots, deren Diskriminierung an der $p$-Identität hängt} -- $(1-0{,}98650)/\xi = 101{,}2$; die $q^{*}$-Bestimmung ist mit der Weg-2-Formel entfallen.)}

\subsection{Zeuge B: die Hochpotenz-Stelle (A270)}

\textbf{[K]} A270 f\"uhrt die Identit\"at
$K_\text{frak}^{-36} \approx 16/\pi^2$ mit dokumentierter
\"Ubereinstimmung $0{,}010\,\%$. Der Exponent $36$ verst\"arkt den
Formabstand auf $36\cdot8{,}9\cdot10^{-5} = 0{,}32\,\%$ -- weit oberhalb
dieser Toleranz. Damit ist die Stelle aufl\"osungsf\"ahig:

\begin{center}
\begin{tabular}{lll}
\toprule
Ausdruck & Wert & rel.\ Abw.\ von $16/\pi^2 = 1{,}6211389$ \\
\midrule
$(74/75)^{-36}$ (additiv) & $1{,}6213007$ & $\mathbf{+1{,}0\cdot10^{-4}}$ \\
$\bigl((1-\xi)^{100}\bigr)^{-36}$ (multipl.) & $1{,}6161261$ & $-3{,}1\cdot10^{-3}$ \\
\bottomrule
\end{tabular}
\end{center}

\textbf{Urteil: 31:1 zugunsten additiv} -- bedingt durch die Annahme, dass
$16/\pi^2$ die richtige Referenz ist. Diese ist in A270 als Koinzidenz
ohne Vorw\"artsherleitung gef\"uhrt (P35), aber Monte-Carlo-abgesichert.

\textbf{[K] Aufwertung der Referenz durch Dok.~314:} Die Konstante ist
korpus-nativ identifizierbar. $\pi^2/16$ ist die
\emph{Packungsdichte des D4-Gitters} (Dok.~314: \enquote{Kusszahl $24$
gegen $8$, Packungsdichte $\pi^2/16$ gegen $\pi^2/32$}) -- also gilt
\[
  \frac{16}{\pi^2} \;=\; \frac{1}{\Delta(\mathrm{D4})}.
\]
Die A270-Identit\"at liest sich damit als
$K_\text{frak}^{-36} = 1/\Delta(\mathrm{D4})$: Der Kehrwert der
Packungsdichte genau des Gitters, dessen Struktur Dok.~314 im
Hilbertraum verifiziert. Die Referenz ist damit keine freischwebende
sch\"one Zahl mehr, sondern eine Gitterinvariante des Rahmens; P35
verengt sich von \enquote{woher die Konstante?} auf \enquote{warum der
Exponent $36$ an die Packungsdichte koppelt}. Die Bedingung von
Zeuge~B wird dadurch deutlich schw\"acher, verschwindet aber nicht.

\subsection{Zeuge C: die Potenzform (A040)}

\textbf{[K]} Die Potenzform $(D_f^\text{eff}/3)^{D_f^\text{eff}/2}$ mit
$D_f^\text{eff} = 2{,}973$ ergibt $0{,}986651$: Residuum additiv
$1{,}6\cdot10^{-5}$, multiplikativ $1{,}0\cdot10^{-4}$ -- tendiert
additiv. Die vierstellige Rundung von $D_f^\text{eff}$
($\pm2{,}5\cdot10^{-4}$ auf $K$) \"ubersteigt jedoch den Formabstand:
kein eigenst\"andiger Zeuge, aber konsistent mit B.

\textit{(Korrigiert am 30.~Sept.~2026: Der Abschnitt \enquote{Nachrechnung der p-Identität} ist entfernt. Er bestimmte den Exponenten $q = -p$ der Weg-2-Formel aus der jeweils angenommenen $K$-Form zurück ($q^{*}_\text{add} = 0{,}267950715$, $q^{*}_\text{mul} = 0{,}267960667$), fand $2-\sqrt3 = 0{,}267949192$ im Abstand $1{,}5\cdot10^{-6}$ als einzigen Kandidaten eines Scans von $175$ Ausdrücken ($\approx10^{-4}$ Zufallswahrscheinlichkeit) und buchte den Rest ($7$~eV in $m_e$) als \enquote{Test der Polmasse}. Die Formel ist aus A130 entfernt; weil sie $m_e/\text{MeV} = 4\sqrt2\,K\,\xi^{q}$ setzt, hängt jede solche Nähe an der Einheit MeV und trägt keine Aussage über die Form von $K_\text{frak}$. Die Skripte \texttt{pruefrechnung\_p\_identitaet.py} und \texttt{pruefrechnung\_rest\_0p1xi.py} sind entsprechend als zurückgezogen markiert.)}

\section{Eingerollt und ausgerollt: warum die additive Form einen Grund hat}
\label{sec:eingerollt}

Abschnitt~\ref{sec:zeugen} behandelt die Formfrage
als Messproblem. Dieser Abschnitt behandelt sie als
\emph{Strukturproblem}: Die beiden Formen sind nicht zwei willk\"urliche
Rechenans\"atze, sondern die exakten Buchhaltungen \emph{zweier
verschiedener Darstellungen derselben Geometrie}. Die Formfrage wird
damit zu einer pr\"azisen physikalischen Frage: In welcher Dom\"ane ist
die Korrektur exakt definiert -- eingerollt oder ausgerollt?

\subsection{Die beiden Dom\"anen}

\textbf{[B]} \emph{Eingerollt (T$^4$, Windungsbild).} Auf dem kompakten
Zyklus ist die nat\"urliche Gr\"o\ss e die Phase bzw.\ Bogenl\"ange, und
Phasen \emph{addieren} sich. Eine Korrektur, die jede Windung um das
Segment $\varepsilon$ verk\"urzt, summiert \"uber $N$ Windungen linear.
Die Schlie\ss ungsbedingung aus A040,
\[
  75\,(1-\varepsilon) = 74 \quad\Rightarrow\quad \varepsilon = \tfrac1{75}
  = 100\xi,
\]
ist in dieser Dom\"ane \emph{exakt} -- als Buchhaltungsidentit\"at der
Windung, nicht als N\"aherung: $75$ verk\"urzte Uml\"aufe schlie\ss en
genau dann, wenn insgesamt ein voller Umlauf fehlt. Die additive Form
$K = 1-100\xi$ ist also keine Linearisierung von irgendetwas; sie ist die
native, geschlossene Buchhaltung der eingerollten Geometrie -- in exakter
Bruchrechnung darstellbar, weil $\xi$ rational ist.

\textbf{[B]} \emph{Ausgerollt (Skalenfluss, Messdom\"ane).} Rollt man den
Zyklus in die Skalenrichtung aus, wird aus der Windungsz\"ahlung ein Lauf
\"uber Gr\"o\ss enordnungen (der RG-Lauf aus A040, $\sim17$ Dekaden). In
dieser Darstellung wirken Korrekturen stufenweise auf Amplituden, und
stufenweise Amplitudenkorrekturen komponieren multiplikativ:
$(1-\xi)^{100}\approx e^{-100\xi}$. W\"are die ausgerollte Beschreibung
fundamental, w\"are \emph{diese} Form exakt und die additive nur ihre
erste Ordnung.

\subsection{Der Formabstand ist der \"Ubersetzungsfehler}

\textbf{[B]} Nach Abschnitt~\ref{sec:formfrage} ist der gesamte
Formabstand der Binomialterm zweiter Ordnung, $4950\xi^2$. Er ist damit
nichts anderes als der Fehler der \"Ubersetzung zwischen den Dom\"anen in
zweiter Ordnung. Die Formfrage lautet scharf: Welche Dom\"ane tr\"agt die
exakte Definition, welche die N\"aherung? Genau das -- und nur das --
messen die Zeugen; aufl\"osungsf\"ahig ist derzeit nur Zeuge~B. \textit{(Korrigiert am 30.~Sept.~2026: vorher \enquote{messen die Zeugen A und B} -- Zeuge~A löst nach Wegfall der Weg-2-Formel nicht auf.)}

\subsection{Die Euler-Parallele macht die Zuordnung eindeutig}

\textbf{[B]} Die Musikspirale liefert den Kontrollfall, in dem die
Dom\"anenzuordnung \emph{unstrittig} ist -- weil dort gemessen wird. Was
Ohr und Messger\"at erreichen, sind Frequenzen, also Gr\"o\ss en der
ausgerollten Skalendom\"ane; Intervalle komponieren dort multiplikativ
($3/2\cdot3/2\cdot\ldots$), und erst der Logarithmus rollt sie auf den
Oktavkreis -- wobei die Irrationalit\"at hereinkommt und die
Schlie\ss ung prinzipiell verhindert. Der $\xi$-Zyklus ist umgekehrt in
der eingerollten Dom\"ane \emph{definiert}: von vornherein eine
Kreisstruktur mit rationaler Windungszahl, ohne zwei multiplikativ
definierte Objekte, die erst logarithmisch vergleichbar gemacht werden
m\"ussten.

\begin{center}
\begin{tabular}{lll}
\toprule
 & Musik (Euler) & FFGFT ($\xi$-Zyklus) \\
\midrule
Messgr\"o\ss e & Frequenz (ausgerollt) & --- (eingerollt definiert) \\
Komposition & Intervalle multiplizieren & Windungsphasen addieren \\
Br\"ucke zum Kreis & Logarithmus (erzwungen) & keine n\"otig \\
Schrittweite & $\log_2(3/2)$ irrational & $1/75$ rational \\
Schlie\ss ung & nie (nur Kommas) & exakt nach 75 \\
\bottomrule
\end{tabular}
\end{center}

Die Zuordnung ist vollst\"andig: multiplikative Komposition geh\"ort zur
ausgerollten Mess-/Skalendom\"ane, additive Komposition zur eingerollten
Windungsdom\"ane.

\subsection{Der Korpus enth\"alt die Gabelung bereits: Fall B gegen Fall C}

\textbf{[K]} Die Zuordnung steht nicht neben dem Korpus, sondern in ihm.
Die Schlie\ss ungsgabelung (Dok.~295/313, zitiert in Dok.~314) f\"uhrt
drei F\"alle: Fall~A ($d=0$, Schlie\ss ung nach einem Umlauf), Fall~B
($\xi$ eingefroren: rationale Rotationszahl $1-d = 74/75$, wegen
$\gcd(74,75)=1$ Schlie\ss ung nach genau $75$ Uml\"aufen) und Fall~C
($\xi$ l\"auft \"uber die Rekursion: \enquote{der kumulierte Defekt wird
logarithmisch und die Windung zur \"aquiangularen Spirale} -- keine
Schlie\ss ung). Das ist exakt die Formfrage als Physik: Fall~B
\emph{ist} die additive Buchhaltung (linearer Defekt, rationale
Windung, exakter Schluss), Fall~C \emph{ist} die
multiplikativ-logarithmische Komposition -- und die \"aquianguläre
Spirale ist dasselbe mathematische Bild wie die Euler-Spirale aus
Kontrollfall~I. Die beiden Formen von $K_\text{frak}$ sind also die
Buchhaltungen der beiden Gabel\"aste; die Frage \enquote{additiv oder
multiplikativ?} ist die Frage, welcher Ast die Korrektur tr\"agt.

Dok.~314 (Kap.~D2) pr\"azisiert zus\"atzlich die Dom\"anengrenze in der
hier verwendeten Terminologie: \emph{Im eingerollten Zustand gibt es
keine fraktalen Defekte; der Defekt $d = 100\xi$ akkumuliert nur
ausgerollt}, und $K_\text{frak}$ ist \emph{eine Eigenschaft der
Durchquerung}, die \"uber die Br\"uckenkonstanten ($E_0$, $v$)
hereinkommt, nicht \"uber den lokalen Operator. Die additive Form ist
danach die Buchhaltung der \emph{geschlossenen} Durchquerung mit
eingefrorenem $\xi$ (Fall~B); die multiplikative w\"are die Buchhaltung
der \emph{laufenden} Akkumulation (Fall~C). Alle
Korpus-Verwendungen von $K_\text{frak}$ (A040-Regel, R70-Zuordnung,
R72) buchen Fall~B.

\subsection{Was die Zeugen dann bedeuten}

\textbf{[B]} In dieser Lesart misst Zeuge~B nicht ein
Rechendetail, sondern die Dom\"anenzugeh\"origkeit der Korrektur -- und
er antwortet: Die exakte Buchhaltung ist die
\emph{eingerollte}; Zeuge~C tendiert in dieselbe Richtung. \textit{(Korrigiert am 30.~Sept.~2026: vorher \enquote{die Zeugen A und B \ldots\ antworten übereinstimmend}.)} Die multiplikative Form ist die N\"aherung, die ein
Beobachter der ausgerollten Skalenerz\"ahlung naiv ansetzen w\"urde --
richtig in erster Ordnung, falsch ab $4950\xi^2$.

Das l\"ost zugleich die scheinbare Spannung innerhalb von A040 auf: Das
Dokument enth\"alt beides -- das Schlie\ss ungskomma (eingerollte, exakte
Lesart) und den kumulativen RG-Lauf (ausgerollte Herleitung des Faktors
$100$). Beide sind konsistent, weil sie in erster Ordnung
\"ubereinstimmen; exakt ist nach Zeugenlage die eingerollte. Der RG-Lauf
leitet her, \emph{warum} die Korrektur die Gr\"o\ss e $100\xi$ hat; die
Windungsgeometrie legt fest, \emph{wie} sie komponiert. Das naive
RG-Argument f\"ur die multiplikative Form wird damit eingeordnet statt
widerlegt: Es ist das korrekte Argument der ausgerollten Dom\"ane -- und
diese ist nach Zeugenlage nicht die definierende.

\subsection{Anschluss an die Grundarchitektur}

\textbf{[B]} Die Zuordnung \enquote{eingerollt $=$ exakt} ist keine
Ad-hoc-Rettung der additiven Form, sondern deckungsgleich mit der
Grundarchitektur des Rahmens: Die kompakte Geometrie (T$^4$, rationale
$\xi$-Setzung) ist prim\"ar; Skalenfluss-Erz\"ahlungen sind abgeleitete
Beschreibungen. Dass die exakte Arithmetik des Korpus durchg\"angig mit
Br\"uchen arbeitet, ist dieselbe Aussage in methodischer Form: Exakte
Br\"uche existieren nur in der eingerollten Dom\"ane -- die ausgerollte
kennt nur N\"aherungen mit $\xi^2$-Resten. Insofern \emph{muss} ein
konsistent eingerollt gebauter Rahmen die additive Form tragen; f\"ande
die Messung multiplikativ, w\"are nicht nur $K_\text{frak}$ betroffen,
sondern die Primatsbehauptung der Geometrie selbst.

\section{Status und Falsifikationslinie}

\begin{keyresult}[Buchungsstand]
$K_\text{frak} = 1-100\xi = 74/75$: Wert \textbf{[K]} (A040). Form
additiv: \textbf{[B]}, einfach bedingt-best\"atigt -- die Bedingung
($16/\pi^2$-Referenz, Zeuge~B) ist durch Dok.~314 wesentlich geschw\"acht, da
$16/\pi^2 = 1/\Delta(\mathrm{D4})$ korpus-nativ identifiziert ist und
nur die Kopplung an den Exponenten~$36$ offen bleibt (P35 verengt).
Strukturell ist die additive Form die Buchhaltung des Gabelasts
Fall~B (eingefrorenes $\xi$, rationale Windung), den alle
Korpus-Verwendungen buchen; die multiplikative Alternative
$(1-\xi)^{100}$ ist die Buchhaltung des Spiral-Asts Fall~C, wird vom
bedingten Zeugen~B disfavorisiert und von keinem Zeugen gest\"utzt.
\end{keyresult}
\textit{(Korrigiert am 30.~Sept.~2026: vorher \enquote{zweifach bedingt-bestätigt}, mit der $p$-Identität $2-\sqrt3$ als Bedingung~1 -- mit der Weg-2-Formel entfallen, siehe Zeuge~A.)}

\textbf{[B]} \emph{Warum trotzdem nicht [K]-gemessen:} Die einzige
scharfe Diskriminierung (Zeuge~B) h\"angt am Status von $16/\pi^2$ als
Referenz (P35). Stellt sich diese Voraussetzung sp\"ater als richtig
heraus, wird die Formaussage unbedingt. Umgekehrt bleibt die
multiplikative Form logisch m\"oglich, solange sie offen ist: Dann w\"are
die $16/\pi^2$-N\"ahe ein Zufall -- und das naive RG-Argument der
ausgerollten Dom\"ane h\"atte recht behalten. \textit{(Korrigiert am 30.~Sept.~2026: vorher zwei Voraussetzungen, die zweite die Herleitung von $p=-(2-\sqrt3)$ für Zeuge~A -- entfallen mit der Weg-2-Formel.)}

\textbf{[B]} \emph{Falsifikationslinie.} Unbedingt entscheidbar w\"are die
Frage an einer Stelle, wo $K$ in hoher Potenz gegen eine \emph{Messung}
statt gegen eine Koinzidenz-Konstante steht. Der Kandidat steht bereits
im Korpus: die Baryon-Vorhersage aus A270 (Sektor-Exponent $+2$, also
$K^{38}$-Niveau, Formabstand dort $\approx0{,}34\,\%$). Eine
Baryon-Massenpr\"ufung mit besserer Genauigkeit entschiede additiv gegen
multiplikativ ohne Bedingungen -- und testete zugleich die
Dom\"anenzuordnung aus Abschnitt~\ref{sec:eingerollt}: Fiele sie
multiplikativ aus, w\"are nicht nur die Formwahl korrigiert, sondern die
Korrektur lebte in der Skalendom\"ane, gegen die Windungslesart des
Schlie\ss ungskommas. Das Strukturargument erh\"oht also den Einsatz; es
bindet die Formfrage an die Primatsfrage.

\section{Offene Punkte}

\begin{itemize}
  \item \textbf{P-315-1:} entfällt. \textit{(Korrigiert am 30.~Sept.~2026: vorher \enquote{Herleitung oder Deklaration der Identität $p=-(2-\sqrt3)$ in A130/Weg~2} -- die Weg-2-Formel ist aus A130 entfernt.)}
  \item \textbf{P-315-2:} Der Rest der $\alpha$-Stelle: $\alpha$ verlangt
    $K_\text{frak} = 0{,}98650$, also $1{,}7\cdot10^{-4}$ neben $74/75$
    (in $n$-Sprache $101{,}2$ statt $100$). Korrekturterm der Theorie oder
    Polmassen-Effekt? Konsistent mit R72. \textit{(Korrigiert am 30.~Sept.~2026: vorher der Rest $1{,}36\cdot10^{-5}$ ($\approx7$~eV in $m_e$) aus der $p$-Identität -- entfallen mit der Weg-2-Formel.)}
  \item \textbf{P35 (verengt durch Dok.~314):} Mit
    $16/\pi^2 = 1/\Delta(\mathrm{D4})$ ist die Konstante als
    D4-Packungsdichte-Kehrwert identifiziert; offen bleibt die
    Vorw\"artsherleitung der Kopplung
    $K_\text{frak}^{-36} = 1/\Delta(\mathrm{D4})$, insbesondere die
    Rolle des Bulk-Exponenten~$36$.
  \item \textbf{P-315-3 ($n$-Sch\"arfe):} In $n$-Sprache erfordert die
    Formdiskriminierung $\Delta n \le 0{,}66$
    ($= n(n-1)\xi/2$). Die drei $n$-Wege aus R72
    ($101{,}2$; $100\pm2$; $100{,}27$) erreichen das nicht; die
    \textit{(Korrigiert am 1.~Okt.~2026: vorher $101{,}3$ -- $(1-\xi m_em_\mu/\alpha)/\xi = (1-0{,}9865006)/\xi = 101{,}246$, wie oben bereits korrigiert.)}
    Sektorleiter ($100{,}27$) ist der aussichtsreichste Kandidat f\"ur
    eine Sch\"arfung.
  \item \textbf{Pr\"azisierung A040:} $D_f^\text{eff}$ mit mehr Stellen
    angeben (ben\"otigt: besser als $\pm1\cdot10^{-4}$), damit Zeuge~C
    aufl\"osungsf\"ahig wird.
  \item \textbf{Baryon-Stelle (A270):} Massenpr\"ufung auf
    $K^{38}$-Niveau als unbedingter Formtest -- zugleich Test der
    Gabelzuordnung Fall~B/C.
\end{itemize}

\section{Verifikation}

Alle Zahlen dieses Dokuments sind durch die ersten beiden Skripte reproduzierbar; die letzten beiden sind zurückgezogen (Weg-2-Formel entfernt) und dokumentieren nur den früheren Stand
(reine Standardbibliothek; exakte Bruchrechnung, wo m\"oglich):

\begin{itemize}
  \item \texttt{python/315\_Skripte/euler\_spirale\_7limit.py} --
    Komma-Suche im 5/7-Limit, exakt; Assertion gegen exakten Schluss.
  \item \texttt{python/315\_Skripte/pruefrechnung\_kfrak\_form.py} --
    Zeugen A/B/C mit Urteilslogik (Zeuge~A jetzt über das von $\alpha$ verlangte $K_\text{frak}$) (Unsicherheit gegen wirksamen,
    ggf.\ verst\"arkten Formabstand).
  \item \texttt{python/315\_Skripte/pruefrechnung\_p\_identitaet.py} --
    exakte $q^{*}$-Bestimmung, Fairness-Scan, Rest- und
    Pr\"afaktor-Gegenprobe.
  \item \texttt{python/315\_Skripte/pruefrechnung\_rest\_0p1xi.py} --
    Korrekturterm-Familien, Kandidatenstatistik
    (Zufallserwartung), Kettenbruch-Diagnose.
\end{itemize}

Eingaben: $\xi=4/30000$ exakt; $m_e = 0{,}51099895069$~MeV,
$m_\mu = 105{,}6583755$~MeV (CODATA~2022).
